
relation est équivalente à la relation :
\[
  f^*\bigl(i_*(\lambda^p(N,y))\bigr)=
  f^*\bigl(\lambda^p(i_*(y))\bigr).
\]
Or on a :
\[
\begin{aligned}
 f^*\bigl(\lambda^p(i_*(y))\bigr)
 &=\lambda^p\bigl(f^*(i_*(y))\bigr)\\
 &=\lambda^p\bigl(j_*(g^*(y)\lambda_{-1}(F))\bigr)
\qquad(3.4)
\end{aligned}
\]
D'autre part :
\[
  f^*\bigl(i_*(\lambda^p(N,y))\bigr)
  =
  j_*\bigl(g^*(\lambda^p(N,y))\lambda_{-1}(F)\bigr)\ (3.4).
\]
Comme \(g^*(N)=F+L\), on peut appliquer la formule (V (5.4.2)) :
\[
  g^*(\lambda^p(N,y))\lambda_{-1}(F)
  =
  \lambda^p(g^*(N),g^*(y))\lambda_{-1}(F)
  =
  \lambda^p(L,g^*(y)\lambda_{-1}(F)).
\]
En posant : \(z=g^*(y)\lambda_{-1}(F)\), on est donc ramené à montrer la relation :
\[
  j_*\bigl(\lambda^p(L,z)\bigr)=\lambda^p(j_*(z)),
\]
c'est-à-dire la relation analogue pour l'immersion \(j\), et l'élément \(z\). Or par hypothèse \(\lambda_{-1}(F)\) est multiple de \((1-L)\); on peut trouver \(z'\) tel que \(z=z'(1-L)\). Alors \(z=j^*(u')\) avec \(u'=j_*(z')\) (2.7) ; comme \(L=j^*(J')\), on a finalement à vérifier :
\[
  j_*\bigl(j^*(\lambda^p(J',u'))\bigr)=
  \lambda^p\bigl(j_*(j^*(u'))\bigr).
\]
La suite exacte
\[
  0\longrightarrow \underline{J}'\longrightarrow \underline{O}_{X'}
  \longrightarrow \underline{O}_{Y'}\longrightarrow 0
\]
donne la relation :
\begin{equation}\tag{4.3.2}
  j_*(1)=1-J'.
\end{equation}
La relation à vérifier s'écrit donc finalement, en appliquant la formule de
projection :
